\section{Discussion}
\label{sec:discussion}
The results distinguish three effects of a mobility budget. Uniform
mobility regularizes each kinetic defect on a local diffusion length.
Predetermined grading changes how the budget is shared among possible
defect positions, delaying the disorder-moment transition from $4/3$
to $8/5$. Observation permits placement at the realized defects and
changes the optimal mean exponent from $1/4$ to $1/5$. The finite-bin
theorem connects these decisions quantitatively: the local uncertainty
must be compared with the placement width, and the resulting conditional
values must still be integrated through the fold configurations.

The strongest coefficients are not all elementary numbers. Above the
predetermined threshold, $\mathcal S_q$ is an attained whole-line
integrated fold-response minimum. At intermediate precision,
$\mathcal F_{\rm ctr}$ is an attained uncertain-center minimum and
$\mathcal H$ is its explicit integrable average. The compactness and
recovery arguments prove exact asymptotic values in both cases.
They require neither an explicit minimizing density nor uniqueness.
For the predetermined supercritical value, the sharp lower bound
requires equal limiting allocations to the two rescaled folds and
rules out wasted limiting mass; achieving only the optimal order
does not impose this allocation.

These distinctions matter when using the results as design guidance.
The mean improvement from predetermined grading is small in its leading
coefficient, and a trial motivated by that limit can perform worse than
uniform mobility at a finite budget. The critical logarithm also
converges slowly. Numerical trial costs therefore cannot be substituted
for the finite-budget optimum, and an asymptotic coefficient does not
by itself prescribe a practical coating at a given physical budget.
The numerical resolution studies in \cref{sec:numerics} make the
separate discretization limitation explicit.

The transport comparison holds at fixed positive transverse bulk
diffusivity, constant positive adsorption affinity, nonzero mean speed,
and a connected one-dimensional wall. It concerns stationary long-time
flow-induced dispersion. Molecular axial diffusion contributes
separately as in \eqref{eq:molecular}; under the stated boundedness
conditions it does not change the divergent leading optima.
The whole-line source models are local energy-dual limits and do not
represent particles in equilibrium on an infinite wall. Likewise,
disorder moments describe a collection of fixed channels, not high
displacement moments or anomalous time dependence within one channel.

The admissible design class idealizes control of tangential diffusivity
independently of affinity, desorption, and geometry. The analysis does
not establish that these controls can be varied independently in a
particular surface chemistry. The theorems concern the specified
integral-budget class; they do not establish optimal values under
additional pointwise, background, or fabrication constraints.
A vanishing bulk diffusivity or a
kinetic family that nearly vanishes everywhere can invalidate the
uniform physical comparison; the Gaussian amplitude-collapse example
shows why local zero geometry alone is insufficient. The generic-fold
theorem requires its stated transverse and sampling assumptions and
does not cover higher degeneracies, simultaneous folds, or unknown
fold locations. Finally, equal-bin observation has no measurement
cost or stochastic noise model. These boundaries specify the transport
and information problem solved by the present laws.
